\documentclass[12pt]{article}
\usepackage[utf8]{inputenc}
\usepackage[T1]{fontenc}
\usepackage{hyperref}
\title{Kadung Language Resources for K-12 Curriculum Development}
\author{Compiled by Claude Sonnet 5.5 (claude-sonnet-5-5) for Andr\'as Kornai}
\date{}
\begin{document}
\maketitle

\section{Introduction}
Kadung (ISO 639-3 \texttt{dkg}; Glottocode vagh1248) is a Plateau (Benue-Congo, Niger-Congo) language of north-central Nigeria, spoken in Plateau and Bauchi states. It is grouped with Ya (Tiyaa) and Bijim under the Kwanka name, with fairly low mutual intelligibility. Speaker estimates are around 19,000--21,000. Very little educational material was found in search; this report is accordingly short.

\section{Primary Educational Resources}
No primers, textbooks or school programmes for Kadung were found in search results.

\section{Linguistic and Reference Materials}
\begin{itemize}
\item Wikipedia, Kwanka language (covers Kadung, Ya and Bijim): \url{https://en.wikipedia.org/wiki/Kwanka_language}
\item ASJP word-list entry for Kadung: \url{https://asjp.clld.org/languages/KADUNG}
\item Glottolog record: \url{https://glottolog.org/resource/languoid/id/vagh1248.bigmap.html}
\end{itemize}

\section{Religious and Community Resources}
\begin{itemize}
\item Scripture Earth page for dkg (search snippets indicate audio recordings from Global Recordings Network and the Nigeria Bible Translation Trust): \url{https://www.scriptureearth.org/00eng.php?iso=dkg}
\item Joshua Project language page for dkg (scripture portions only): \url{https://joshuaproject.net/index.php/languages/dkg}
\item Joshua Project people profile, Kadun in Nigeria: \url{https://joshuaproject.net/people_groups/12912}
\end{itemize}

\section{Digital and Audio Resources}
Only the audio recordings listed on the Scripture Earth page above were indicated; their contents were not verified.

\section{Curriculum Development Context}
No evidence of Kadung in schools was found. A sociolinguistic profile by Decker et al.\ (2021) on Bijim, mentioning Kadung and Tiyaa, was referred to in search summaries, but no URL for it was seen, so it is not cited. A curriculum would need to begin with orthography and primer development in consultation with the community.

\end{document}